\section{Particularidades de cada instalaci\'on}

Una vez entendido el funcionamiento general, es mejor consultar cada instalaci\'on en una unidad separada.

Este reparto reduce el tama\~no del manual principal, facilita la revision por especialidad y evita que una correccion local obligue a navegar por un documento demasiado largo.

\subsection{Unidades disponibles}

\begin{itemize}
\item \manualexterno{MANUAL_USUARIO_AGUA.pdf}{Instalaci\'on de Agua}
\item \manualexterno{MANUAL_USUARIO_SANEAMIENTO.pdf}{Instalaci\'on de Saneamiento}
\item \manualexterno{MANUAL_USUARIO_VENTILACION.pdf}{Instalaci\'on de Ventilaci\'on}
\item \manualexterno{MANUAL_USUARIO_ELECTRICIDAD.pdf}{Instalaci\'on de Electricidad}
\end{itemize}

\subsection{Como se relacionan los documentos}

La logica propuesta es esta:

\begin{enumerate}
\item El manual global reune el flujo comun a todas las instalaciones.
\item Cada instalaci\'on mantiene su unidad propia con sus reglas, avisos y capturas especificas.
\item Los enlaces entre PDFs permiten saltar desde el indice global a la unidad que corresponda y volver despues al documento principal.
\end{enumerate}

\subsection{Criterio editorial}

Cuando una regla afecte a todas las instalaciones, debe mantenerse en el manual global.

Cuando una regla dependa del sistema, del catalogo de piezas, de los diametros, de los layers o de avisos especificos, debe ir a la unidad de la instalaci\'on correspondiente.
